\properlane{turned-to-god}{The people who ask God to turn them}

\subsection{Can a confession contain an accusation?}

Israel wants God to find a people who remember him and do right. A few lines later the same people say that none of them calls on his name. If they cannot even bring themselves to pray, who is speaking the words Isaiah has given them? The problem begins with what they say to God about their wandering. After calling him their father they ask why he allowed them to get so far from him. In the language Jerome knew, they ask God why he \emph{made} them stray. One could hear a person owning up to a fault and then laying the responsibility for it elsewhere.

Jerome starts elsewhere in Scripture, with a letter written to an entire city church. Paul greets the church at Corinth along with those everywhere who call on the Lord (1 Cor 1:2). Within the Corinthian church itself there are married people, people considering whether to marry, people who have let themselves be drawn into idolatry and people who need instruction about the resurrection. Paul can praise the Corinthians in one place and correct them in another. Nobody need infer that every Corinthian is equally praiseworthy on every count or guilty of each fault. A people at prayer can likewise include saints and sinners. Some of their words give God praise. Other words from the same gathering contain a sinner's attempt to account for his wrong conduct by questioning God. Jerome takes the prayer seriously enough to distinguish the speakers whose voices it lets us hear.\footnote{Jerome, \textit{Commentary on Isaiah}, book XVII, on Is 63:16--19, Latin part IV, pp.~149--151.}

Jerome also attends to a genuine experience behind the charge. A man does wrong and God appears to let it go. The man goes further and the habit hardens. If he later asks how God could have let this happen, he remembers an actual period during which he was free to persist. What he could not see at the time was the intention of the one who waited for him. Jerome compares a doctor with a frightened patient. The doctor has a sharp instrument. If he displays it before he can treat the man, the sight may terrify his patient. He keeps it out of sight. For Jerome the sinner mistakes God's patience for permission to go on. The delay that left time for him to turn has allowed him to tell himself there will never be anything to answer for.

At the end of the prayer the worshippers speak of clay and call the one who made them their Father. Whatever clay becomes is owing to the hands that work it. A vessel has no independent claim against its maker on the ground that it might once have made a better job of itself. Yet these speakers can address their maker as Father. That difference is where Jerome locates their daring. Children have a reason to speak to their father even while accepting his correction. All the excuses with which the community might have made itself out to be innocent can fall away and it can still pray. What it asks the Father to remember is the identity of those who now bear the consequences of their acts.\footnote{Jerome, \textit{Commentary on Isaiah}, XVII, on Is 64:7--8, pp.~158--159.}

The wording of the two Isaiahs a reader may encounter makes a difference here. In Jerome's Latin text, reflected in Challoner 64:5, God has met a just man and the people who confess their sin go on to say they will be saved. Jerome asks how such people can speak with assurance. They have been in sin a long time. The only source of their confidence that remains is God's mercy. At Mass in the United States the speakers wish God might meet \emph{them} doing right. Then come the admission of sin and the vivid description of their own uncleanness. A listener following the proclamation hears their wish tested almost before they have finished expressing it. Can anyone bring the deeds asked for in the Collect if even Israel's acts of justice have been stained?\footnote{Jerome, \textit{Commentary on Isaiah}, XVII, p.~155, on the clause following the admission of sin in his text of Is 64.}

Thomas Aquinas puts that general accusation into terms a person examining an actual action can understand. There is the work undertaken and then performed carelessly. There is a sound deed chosen with a crooked intention. There is pride that arises even as someone recognizes a thing well done. A person charged with distributing money for food can injure the people it was meant to reach by neglecting to check where the money goes. A person who attends diligently to its use might still use the act to secure praise. And a person who originally helped for the neighbor's sake might take complete credit once the good result became known. These failures call for different corrections. The task is worth doing on each account. Asking God for the will to do it well includes a petition concerning what happens in the heart before, during and after the work.\footnote{Thomas Aquinas, \textit{Expositio super Isaiam ad litteram}, chapter 64, §85024, on the ways acts of justice can be stained.}

\subsection{Who makes the first turn?}

Jerome turns to the Psalter as he explains why Isaiah's people can still pray when God seems to hide his face. He pictures them confessing their sin, hearing God's promise of pardon and then praying the words of Psalm 79:4 for the sight of God's face and for salvation.\footnote{Jerome, \textit{Commentary on Isaiah}, XVII, p.~158, on Is 64:7.} The congregation sings that petition in response to Isaiah at Mass. In Augustine's Psalter the imperative in the response is addressed to God. The people are asking him to turn them. Elsewhere in the psalm they ask \emph{him} to turn again and look down upon the vineyard. A change of direction can be requested on both sides of this damaged relation. The two requests make different demands. The land is already God's and he has already planted the vine. The vine cannot undertake God's work of caring for it. When the people ask for the plant to be brought to completion they appeal to the planter to finish what he began. Augustine can therefore hear in their prayer a confession about the capacity of those who pray. They depend upon the one they address for the very conversion they desire.\footnote{Augustine, \textit{Exposition on Psalm} 79, §§3, 9, 11, on the refrain, the planted vine and the final promise.}

The people ask for life before they promise to call on God. Augustine connects the two. Someone made alive can invoke the one who gave him life. Here is a way of hearing Isaiah's people when their confession has run out of excuses. Asking to be restored already involves the smallest beginning of the turn for which they pray. It can come from God as surely as the further acts they wish to perform. At the Entrance the Church has already begun to lift her soul from what once dragged it down. At the Collect she asks for the willing heart with which to go out to meet Christ. At the psalm she can acknowledge the one to whom she owes both the first motion and its continuation.

Such a person still needs the warning in Mark. Peter promises that others may fail but he will remain faithful. Before the night is out Jesus addresses him by name when he finds him sleeping (Mark 14:29--42). The confidence of the earlier promise did little for the friend who wanted someone awake beside him. In that garden Jesus specifically asks his disciples to pray against temptation. Even in an account in which the disciples have heard and understood that an hour of danger has come, they still need strength to respond in that hour. The congregation that has just asked for mercy before the Gospel is in a position to understand the discrepancy. The Corinthians too could possess real gifts of knowledge and speech and still need correcting. Their gifts gave them the means to attend to Paul when he wrote. Listening to him would require them to change what they were doing.

\subsection{How a person's pleasure can change}

At Communion Psalm 85 speaks of fruit from the earth. The preceding verse speaks of truth springing up from the earth. Augustine gives more than one account of that earlier verse. Christ is Truth born of a woman. He also understands the sinner who finally speaks truth about himself as someone in whom truth has sprung up. Where a person has spent years explaining how everyone else is to blame, simply saying what he himself did amounts to new growth. Augustine compares the condition of the hearer who confesses to the condition of a field broken up by the plough. There is now somewhere for the seed of a word to go. The preacher can sow when he speaks and the hearer can pay attention. Neither can supply the rain that will be needed after they have parted.\footnote{Augustine, \textit{Exposition on Psalm} 84, §§10--12. The truth springing from earth is Challoner Ps 84:12, the fruit yielded by the earth is v.~13.}

At the actual fruit verse Augustine notices a word for sweetness or delight in his version of the psalm. He asks where delight in justice comes from. Augustine describes a thief who took pleasure in stealing what someone else possessed and he lacked. Then he finds himself with something another person lacks and enjoys giving it away. The change goes further than a fear of being caught, or even the moment at which a stolen object is restored to its owner. The thief had experienced someone else's possession as an occasion for acquisition. Now another person's need gives him pleasure in an act that diminishes his own possessions. Augustine attributes this sweetness to God. It is recognized in a use of the hands which would once have seemed to this same person an intolerable loss.

The congregation hearing the Communion antiphon has already offered to God things that came from him. On Augustine's understanding of the psalm we may see another instance of such a gift in the willingness to part with goods for a neighbor. No one can bring that desire to God as an achievement that puts him in their debt. Yet there is a neighbor who actually receives a gift, someone whose hunger can be relieved. In praying after Communion for a deeper love of what endures, communicants who have to handle things that pass away are asking for help with this very sort of decision. A person can learn to part with money precisely because he has learned to love someone whom that money can help.

The face the people ask to see in the response is the face of the one whose humble coming the priest recalls in the Preface. When grace has already reached people who did not know how to return, they have reason to hope that they can come before him openly when he appears. Those who receive the Advent solemn blessing can say Amen to a petition for charity in the present race even if, at the beginning of Mass, they had to own up to its absence. They have heard why God can be asked for something which is also truly theirs to do.

\begin{fourSenses}
\item[Literal.] Israel calls upon its Father from a condition of guilt it can neither conceal nor repair alone. The vine needs its planter's attention. The Corinthians still need correction after receiving gifts for which Paul gives thanks.
\item[Allegorical.] Augustine hears Christ speaking in his Church when the psalmist lifts his soul to God. He recognizes Christ in the salvation for which the people ask and in the face whose glory was hidden in the humiliation of the cross.
\item[Moral.] Conversion reaches what a person takes pleasure in doing. Augustine's former thief can now see a lack in another person as a reason to give. A real good passes between their hands.
\item[Anagogical.] Those who have learned to ask for mercy can hope to look on Christ in glory. The people who ask for a place at his right hand entrust themselves to the one whose gift first enabled them to seek his face.
\end{fourSenses}
